% ==========================================================================
% Chapter 4 — LLM Dynamic Implementation: The "Mixing Board" Approach
% ==========================================================================
\section{LLM Dynamic Implementation}
\label{ch:dynamic_implementation}

This chapter details the first integration strategy: LLM-driven dynamic reward shaping,
informally referred to as the \emph{Mixing Board} approach.
The LLM does not act inside the environment; instead, it periodically
adjusts a vector of subtask priority weights that govern
(i)~which navigation goal the DAG planner selects and
(ii)~the magnitude of the potential-based shaping signal each agent receives.

The exposition follows the exact dataflow of the implementation:
\textbf{DAG-conditional planning} (\cref{sec:dag_planner}) $\rightarrow$
\textbf{LLM weight query} (\cref{sec:llm_weight_query}) $\rightarrow$
\textbf{softmax budget normalisation} (\cref{sec:softmax_budget}) $\rightarrow$
\textbf{potential-based reward shaping} (\cref{sec:pbrs}).

% ──────────────────────────────────────────────────────────────────────────
\subsection{System Overview}
\label{sec:dyn_overview}

\Cref{fig:dyn_system_flow} illustrates the end-to-end pipeline.
At every training update the critic trigger module
(\cref{sec:critic_trigger}) monitors TD-error statistics and
success-rate trends. When a learning plateau is detected, a structured
prompt is constructed, sent to the LLM, and the returned weight vector is
processed by the \texttt{RewardInjector} safety layer before being
applied to the DAG planner and the shaping function.

\begin{figure}[H]
\centering
\begin{tikzpicture}[
    box/.style={draw, rounded corners, minimum width=2.8cm, minimum height=0.9cm,
                font=\small\sffamily, align=center, fill=gray!8},
    arr/.style={-{Stealth[length=2.5mm]}, thick},
    node distance=1.6cm and 2.2cm
]
    % Row 1
    \node[box] (env)    {Env Step\\$s_t \to s_{t+1}$};
    \node[box, right=of env] (planner) {DAG Planner\\$\pi_{\text{goal}}(\boldsymbol{w})$};
    \node[box, right=of planner] (shaping) {PBRS\\$F(s,s',g)$};

    % Row 2
    \node[box, below=1.8cm of env] (trigger)  {Critic Trigger\\(Plateau + Z-score)};
    \node[box, right=of trigger] (llm)     {LLM Query\\(Logit output)};
    \node[box, right=of llm]     (injector){RewardInjector\\(Softmax + EMA)};

    % Arrows Row 1
    \draw[arr] (env) -- node[above,font=\footnotesize]{inventory} (planner);
    \draw[arr] (planner) -- node[above,font=\footnotesize]{$g_t^{(i)}$} (shaping);
    \draw[arr] (shaping.north) -- ++(0,0.5) -| node[near start,above,font=\footnotesize]{$r_t^{\text{shaped}}$} (env.north);

    % Arrows Row 2
    \draw[arr] (trigger) -- node[above,font=\footnotesize]{prompt} (llm);
    \draw[arr] (llm)     -- node[above,font=\footnotesize]{$\boldsymbol{c}^{\text{raw}}$} (injector);
    \draw[arr] (injector.north) -- ++(0,0.5) -- node[above,font=\footnotesize]{$\boldsymbol{w}$} (planner.south);

    % Vertical feed
    \draw[arr,dashed] (env.south) -- node[left,font=\footnotesize]{TD stats} (trigger.north);
\end{tikzpicture}
\caption{Dataflow of the LLM Dynamic ``Mixing Board'' approach.
         Solid arrows denote per-step operations; the dashed arrow
         indicates per-update monitoring.}
\label{fig:dyn_system_flow}
\end{figure}


% ──────────────────────────────────────────────────────────────────────────
\subsection{DAG-Conditional Planner}
\label{sec:dag_planner}

The crafting environment enforces a strict dependency graph over subtasks:
\[
\underbrace{(\text{Wood} + \text{Stone})}_{\text{collect}}
\;\longrightarrow\;
\text{Pickaxe}
\;\longrightarrow\;
\text{Iron}
\;\longrightarrow\;
\underbrace{(\text{Sword} + \text{Armor})}_{\text{craft}}
\;\longrightarrow\;
\text{Bridge}
\;\longrightarrow\;
\text{Enemy}
\;\longrightarrow\;
\text{Gold}
\]

\noindent
\textbf{Goal selection.}
For each parallel environment, the planner inspects the current
inventory vector $\mathbf{x}\in\mathbb{Z}^{10}$ and builds a set of
\emph{feasible} subtasks for each agent.
Agent~0 (Lumberjack) may target \{Wood, Workbench, Bridge, Enemy\};
Agent~1 (Miner) may target \{Stone, Workbench, Iron, Enemy, Gold\}.
A subtask is feasible only if its DAG preconditions are met (e.g.\
\texttt{Iron} requires $\mathtt{pickaxe}\ge 1$).

For each feasible subtask $j$, a \emph{score} is computed:
\begin{equation}
\label{eq:greedy_score}
    s_j = R_j^{\text{base}} \cdot w_j
\end{equation}
where $R_j^{\text{base}}$ is a hand-crafted base reward reflecting the
task's importance in the DAG (e.g.\ $R_{\text{gold}}=15$,
$R_{\text{enemy}}=10$, $R_{\text{bridge}}=3$), and $w_j$ is the
LLM-provided weight for subtask $j$.
The agent greedily selects the subtask with the highest score:
\begin{equation}
\label{eq:greedy_select}
    g^{(i)} = \arg\max_{j \in \mathcal{F}^{(i)}(\mathbf{x})}\;
              R_j^{\text{base}} \cdot w_j
\end{equation}
where $\mathcal{F}^{(i)}(\mathbf{x})$ denotes the feasible set for
agent~$i$ given inventory~$\mathbf{x}$.

The following code excerpt shows the core of this selection for Agent~0:
\begin{lstlisting}[style=code, language=Python, caption={Greedy score computation for Agent~0 in \texttt{orchestrator.py}.}, label={lst:greedy_planner}]
# Agent 0 (Lumberjack) valid tasks
a0_scores = {}
if w < 2:
    a0_scores["Wood"] = R_WOOD * w_wood
needs_wb_a0 = (p == 0) or (sw == 0) or (a == 0)
if needs_wb_a0:
    a0_scores["Workbench"] = R_WORKBENCH * w_workbench
if b == 0:
    a0_scores["Bridge"] = R_BRIDGE * w_bridge
if e == 0:
    a0_scores["Enemy"] = R_ENEMY * w_enemy

a0_target = max(a0_scores.items(),
                key=lambda x: x[1])[0] \
            if a0_scores else "None"
\end{lstlisting}

\noindent
\textbf{Cooperative coupling guardrail.}
If either agent selects a cooperative task (Workbench or Bridge) that
requires both agents to be co-located, the planner pulls the
other agent to the same target:
\begin{lstlisting}[style=code, language=Python, caption={Cooperative coupling in \texttt{orchestrator.py}.}, label={lst:coop_coupling}]
if a0_target == "Workbench" and needs_wb_a1:
    a1_target = "Workbench"
elif a1_target == "Workbench" and needs_wb_a0:
    a0_target = "Workbench"
elif a0_target == "Bridge" and "Bridge" in a1_scores:
    a1_target = "Bridge"
\end{lstlisting}

\noindent
\textbf{DAG guardrails.}
Before any weight vector is used, a hard programmatic filter
(\texttt{apply\_dag\_guardrails}) zeroes out weights whose
preconditions are provably unmet:
\begin{lstlisting}[style=code, language=Python, caption={DAG guardrail enforcement in \texttt{orchestrator.py}.}, label={lst:dag_guardrail}]
def apply_dag_guardrails(weights, metrics):
    guarded = dict(weights)
    # Iron requires Pickaxe
    if metrics.get('pickaxe', 0.0) < 1.0:
        guarded['w_iron'] = 0.0
    # Enemy and Gold require Bridge+Sword+Armor
    if (metrics.get('bridge', 0.0) < 1.0
        or metrics.get('sword', 0.0) < 1.0
        or metrics.get('armor', 0.0) < 1.0):
        guarded['w_enemy'] = 0.0
        guarded['w_gold'] = 0.0
    return guarded
\end{lstlisting}

This is essential because the LLM occasionally hallucinated positive
weights for unreachable subtasks (e.g.\ assigning $w_{\text{iron}}>0$
when no pickaxe exists), injecting noise into the value network.
The programmatic guardrail acts as a hard constraint layer that
overrides the LLM output.

% ──────────────────────────────────────────────────────────────────────────
\subsection{Two-Stage Critic Trigger}
\label{sec:critic_trigger}

The LLM is not queried every update; instead, a \emph{two-stage gate}
determines when an intervention is necessary.

\subsubsection{Stage~1: Plateau Gate}

A rolling window of the last $W$ success-rate values is maintained.
The gate opens when the improvement over the previous window falls
below a threshold~$\epsilon$:
\begin{equation}
\label{eq:plateau_gate}
    \text{gate\_open}
    \;\Longleftrightarrow\;
    \left|\;
        \underbrace{\frac{1}{W}\sum_{k=0}^{W-1}
            \rho_{t-k}}_{\bar{\rho}_{\text{recent}}}
        \;-\;
        \underbrace{\frac{1}{W}\sum_{k=W}^{2W-1}
            \rho_{t-k}}_{\bar{\rho}_{\text{older}}}
    \;\right|
    \;<\;
    \epsilon
\end{equation}
where $\rho_t$ is the success rate at update~$t$, $W$ is the
\texttt{plateau\_window} (default 10), and $\epsilon$ is the
\texttt{plateau\_threshold} (default 0.02).
Additionally, the gate opens if any individual subtask
has plateaued (variance $< 0.02$ over 20~updates while
remaining between 5\% and 95\% completion).

\subsubsection{Stage~2: Z-Score Trigger}

Once the gate is open, the system checks whether the critic's
TD-error variance has spiked beyond $\sigma$ standard deviations from
its running mean, indicating that the value network is ``confused'':
\begin{equation}
\label{eq:zscore_trigger}
    z = \frac{
        \text{Var}[\delta_t] - \overline{\text{Var}[\delta]}
    }{
        \text{Std}\bigl[\text{Var}[\delta]\bigr]
    }
    \;>\;
    \frac{\sigma_{\text{base}}}{\max(\eta,\, 0.1)}
\end{equation}
where $\delta_t$ is the TD error at update~$t$, the bar denotes the
running mean over a window of 100 updates, $\sigma_{\text{base}}=1.5$,
and $\eta$ is a sensitivity factor that decays multiplicatively by
$0.995$ after each trigger, raising the effective threshold over time.

If the gate is open and the trigger fires, an LLM intervention is dispatched.
A cooldown of 25~updates prevents repeated firing.

% ──────────────────────────────────────────────────────────────────────────
\subsection{LLM Weight Query: Continuous Logit Output}
\label{sec:llm_weight_query}

When the trigger fires, the \texttt{PromptBuilder} constructs a
structured prompt containing:

\begin{enumerate}[nosep]
    \item The trigger reason (plateau description, z-score value).
    \item Current subtask completion percentages, with saturation flags ($\ge 95\%$).
    \item Per-agent TD-error statistics (mean, variance, per-agent breakdown).
    \item Current reward shaping weights.
    \item An episodic memory buffer of the last 3 interventions and their
          resulting TD-error deltas.
\end{enumerate}

The prompt explicitly instructs the LLM to output \emph{unnormalised
log-probabilities (logits)}, not final weights:

\begin{quote}
\small\sffamily
``Your output values represent unnormalized log-probabilities (logits).
They will be processed via a temperature-scaled softmax function and
scaled to maintain a constant total reward budget.
Focus on the RELATIVE MAGNITUDE and PRIORITY RANKING of subtasks.''
\end{quote}

This design decouples the LLM's task from absolute magnitude calibration.
The LLM is free to express priority differences via wide logit ranges
(e.g.\ a bottleneck at 8.0 vs.\ saturated at 0.5), and the downstream
normalisation ensures a bounded, stable reward budget.

The LLM responds with a JSON object:
\begin{lstlisting}[style=code, language=Python, caption={Schema of the LLM response.}, label={lst:llm_response}]
{
  "reasoning": "<1-2 sentences>",
  "w_wood":      <float, logit>,
  "w_stone":     <float, logit>,
  "w_workbench": <float, logit>,
  "w_iron":      <float, logit>,
  "w_bridge":    <float, logit>,
  "w_enemy":     <float, logit>,
  "w_gold":      <float, logit>
}
\end{lstlisting}

% ──────────────────────────────────────────────────────────────────────────
\subsection{Softmax Attention Budget}
\label{sec:softmax_budget}

The raw logits $\boldsymbol{c}=(c_1,\dots,c_K)$ returned by the LLM are
converted into a \emph{budget-preserving} weight vector via a
temperature-scaled softmax:
\begin{equation}
\label{eq:softmax_budget}
\boxed{
    w_i
    \;=\;
    K \;\cdot\;
    \frac{
        \exp\!\bigl(c_i / \tau\bigr)
    }{
        \displaystyle\sum_{j=1}^{K}
        \exp\!\bigl(c_j / \tau\bigr)
    }
}
\end{equation}
where $K=|\mathcal{W}|=7$ is the number of subtask weights and
$\tau=0.5$ is the temperature hyperparameter.

\noindent
\textbf{Properties of this equation:}
\begin{itemize}[nosep]
    \item \emph{Constant budget}: $\sum_{i=1}^{K} w_i = K$ always.
          The total reward magnitude is invariant to the LLM's output scale,
          preventing the LLM from accidentally inflating or deflating
          the overall reward.
    \item \emph{Sharpening via temperature}: With $\tau=0.5 < 1$, the
          softmax amplifies differences between logits, producing
          more decisive priority assignments.
          Conversely, $\tau\to\infty$ yields uniform weights.
    \item \emph{Numerical stability}: The implementation shifts logits
          before exponentiation:
\end{itemize}

\begin{lstlisting}[style=code, language=Python, caption={Softmax attention budget in \texttt{reward\_injector.py}.}, label={lst:softmax_code}]
logits = np.array([clipped[k] for k in keys])
logits_shifted = logits - np.max(logits)        # stability
exp_logits = np.exp(logits_shifted / self.softmax_temperature)
softmax_probs = exp_logits / np.sum(exp_logits)
scaled = softmax_probs * len(keys)              # budget = K
\end{lstlisting}

\subsubsection{EMA Smoothing}
\label{sec:ema_smoothing}

To avoid abrupt weight jumps that would destabilise PPO's value function,
the softmax output is blended with the previous weights via
Exponential Moving Average:
\begin{equation}
\label{eq:ema}
    \widetilde{w}_i^{(t)}
    \;=\;
    \alpha \cdot w_i^{\text{softmax}}
    + (1-\alpha) \cdot \widetilde{w}_i^{(t-1)}
\end{equation}
with $\alpha=0.2$. This means that each LLM intervention takes
roughly $\lceil 1/\alpha \rceil = 5$ updates to fully propagate.

\subsubsection{Safety Bounds}
\label{sec:safety}

Before softmax normalisation, individual logits are clipped to
$[0, w_{\max}]$ with $w_{\max}=3.0$, and the raw sum is bounded
by $W_{\text{bound}}=15.0$:
\begin{equation}
    c_i' = \text{clip}(c_i, 0, w_{\max}),
    \qquad
    \text{if}\; \sum_i c_i' > W_{\text{bound}}
    \;\Rightarrow\;
    c_i'' = c_i' \cdot \frac{W_{\text{bound}}}{\sum_j c_j'}
\end{equation}

These bounds prevent adversarial or malformed LLM outputs from
creating extreme weight vectors.

\subsubsection{Auto-Rollback}
\label{sec:rollback}

After each intervention, the \texttt{RewardInjector} monitors the
average environment reward for \texttt{rollback\_window}=25~updates.
If the post-intervention average drops more than 20\% below the
pre-intervention level, the weights are automatically rolled back:
\begin{equation}
    \text{rollback}
    \;\Longleftrightarrow\;
    \frac{
        \bar{r}_{\text{pre}} - \bar{r}_{\text{post}}
    }{
        |\bar{r}_{\text{pre}}|
    }
    > 0.20
\end{equation}


% ──────────────────────────────────────────────────────────────────────────
\subsection{Potential-Based Reward Shaping}
\label{sec:pbrs}

Once the planner has assigned goal zone $g^{(i)}_t$ to agent~$i$,
a shaping reward is computed via the potential-based reward shaping
(PBRS) framework~\cite{Ng_Harada_Russell_1999}. The key theoretical guarantee of PBRS
is that it preserves the set of optimal policies: adding $F$ to
the reward does not change which policy is optimal.

\subsubsection{Potential Function}
The potential of agent~$i$ in state~$s$ with respect to goal zone~$g$ is:
\begin{equation}
\label{eq:potential}
    \Phi^{(i)}(s, g) = D_{\max} - \|\mathbf{p}^{(i)}(s) - \mathbf{z}_g\|_2
\end{equation}
where $\mathbf{p}^{(i)}(s)\in\mathbb{R}^2$ is agent~$i$'s position,
$\mathbf{z}_g\in\mathbb{R}^2$ is the target zone's position, and
$D_{\max}=85$ is an upper bound on the grid diagonal.
The potential is high when the agent is close to the goal and low when far.

\subsubsection{Shaping Reward}
The per-step shaping reward follows the standard PBRS difference form:
\begin{equation}
\label{eq:shaping_reward}
    F_t^{(i)}
    \;=\;
    \bigl(\gamma\,\Phi^{(i)}(s_{t+1}, g) - \Phi^{(i)}(s_t, g)\bigr)
    \;\cdot\;
    \lambda_{\text{scale}}
    \;\cdot\;
    \mathbb{1}[g \text{ active}]
\end{equation}
where $\gamma=0.99$ is the discount factor and
$\lambda_{\text{scale}} = 0.05$ ($\texttt{R\_LLM\_SCALE}$).

The vectorised implementation processes all environments and agents in parallel:

\begin{lstlisting}[style=code, language=Python, caption={Vectorised PBRS in \texttt{orchestrator.py}.}, label={lst:pbrs_code}]
dist_next = np.linalg.norm(pos_next - goal_targets, axis=2)
dist_prev = np.linalg.norm(pos_prev - goal_targets, axis=2)

MAX_DIST = 85.0
phi_next = MAX_DIST - dist_next
phi_prev = MAX_DIST - dist_prev
F = (gamma * phi_next - phi_prev) * R_LLM_SCALE * goal_active
\end{lstlisting}

\subsubsection{Design Decision: No Per-Weight Scaling of $F$}
\label{sec:no_f_scaling}

An important design decision, documented in the codebase, is that the
shaping signal $F$ is \textbf{not} multiplied by the per-subtask weight
$w_j$. The LLM's influence is already applied at goal \emph{selection}
time (\cref{eq:greedy_select}). Scaling the shaping gradient by
transient EMA-smoothed weights would create a non-stationary reward
function whose variance confuses PPO's value network, leading to
training instability.

\subsubsection{Shaping Decay ($\lambda$)}
The shaped reward is blended with the environment reward using a
time-varying coefficient $\lambda_t$:
\begin{equation}
\label{eq:total_reward}
    r_t^{\text{total}}
    = r_t^{\text{env}}
    + \lambda_t \cdot F_t
\end{equation}
$\lambda_t$ begins at 1.0 and anneals linearly to 0.01 once the
bridge-building success rate exceeds 40\%, over 500~updates:
\begin{equation}
    \lambda_t =
    \begin{cases}
        1.0 & \text{if bridge\_pct} < 0.4 \\[4pt]
        1.0 + (0.01 - 1.0) \cdot \min\!\left(\dfrac{d}{500},\, 1\right)
        & \text{otherwise}
    \end{cases}
\end{equation}
where $d$ counts updates since the decay was activated.
This schedule ensures dense guidance in early training (when the agents
have not yet learned the task structure) and a graceful transition to
the sparse environment reward once agents have reached the critical
``bridge'' milestone.

% ──────────────────────────────────────────────────────────────────────────
\subsection{Async Communication Architecture}
\label{sec:async_bridge}

To avoid blocking the training loop during LLM inference, queries are
dispatched through the \texttt{LLMBridge} module, which runs a
background worker thread.

\begin{enumerate}[nosep]
    \item The critic trigger fires and constructs a prompt.
    \item The prompt is enqueued to a thread-safe \texttt{queue.Queue}.
    \item A daemon worker thread dequeues and calls the LLM backend
          (Ollama HTTP API, HuggingFace local, or Gemini API).
    \item On completion, the callback parses the JSON, passes it to
          the \texttt{RewardInjector}, and stores the result as
          ``pending weights.''
    \item At the \emph{next} training update, the
          \texttt{CriticTrigger.on\_update\_end} detects the pending
          weights and applies them.
\end{enumerate}

This design guarantees that the PPO update is never blocked by network
latency. The LLM bridge also maintains a thread-safe response cache,
avoiding redundant queries for identical prompts.
Retry logic uses exponential backoff ($2^k$ seconds for attempt~$k$,
up to 3~retries).

% ──────────────────────────────────────────────────────────────────────────
\subsection{Summary}
\label{sec:dyn_summary}

The Mixing Board approach treats the LLM as a \emph{meta-level supervisor}
that reads aggregate training statistics and outputs priority logits.
A cascade of safety layers (DAG guardrails, clipping, softmax budget,
EMA smoothing, auto-rollback) ensures that the LLM's influence is
bounded and stable, while the PBRS framework guarantees policy-invariance
of the shaping signal.

\begin{table}[H]
\centering
\caption{Key hyperparameters of the Mixing Board approach.}
\label{tab:dyn_hyperparams}
\begin{tabular}{@{}lll@{}}
\toprule
\textbf{Component} & \textbf{Parameter} & \textbf{Value} \\
\midrule
Softmax budget   & Temperature $\tau$            & 0.5 \\
                 & Budget size $K$               & 7   \\
EMA smoothing    & $\alpha$                      & 0.2 \\
Safety bounds    & $w_{\max}$ (per-weight clip)   & 3.0 \\
                 & $W_{\text{bound}}$ (sum clip)  & 15.0 \\
Critic trigger   & Plateau window $W$            & 10  \\
                 & Plateau threshold $\epsilon$  & 0.02 \\
                 & Z-score $\sigma_{\text{base}}$ & 1.5 \\
                 & Cooldown                      & 25 updates \\
                 & Sensitivity decay             & 0.995 \\
Auto-rollback    & Window                        & 25 updates \\
                 & Drop threshold                & 20\% \\
PBRS             & Scale $\lambda_{\text{scale}}$ & 0.05 \\
                 & $D_{\max}$                    & 85.0 \\
Shaping decay    & $\lambda_{\text{init}}$       & 1.0 \\
                 & $\lambda_{\text{final}}$      & 0.01 \\
                 & Decay trigger                 & bridge $>$ 40\% \\
                 & Decay duration                & 500 updates \\
\bottomrule
\end{tabular}
\end{table}
